\begin{abstract}
Small self-play runs, of the kind used in teaching, hobby engines and research prototypes, cannot afford
large evaluations, so progress is often tracked by matching each new network against its predecessor and
summing the Elo differences (``chained Elo'').
We train an AlphaZero-style 9$\times$9 Go engine from scratch on one 8-core laptop CPU
(\valNGens{} generations, \valTotalGames{} self-play games) and report a case where this signal pointed
the wrong way: chained Elo fell after generation~\valChainedPeakGen{} and ranked generation~34 below
generations~19 and~25, yet in direct matches generation~34 beat them in \valHTtfVsNt{} and \valHTtfVsTf{}
games.
Anchored matches against fixed earlier and later generations (\valHtoHGames{} games) show steady
improvement throughout.
Match noise contributes, but cannot explain the size of the error. Replaying the same consecutive
matches with more search per move recovers the correct direction and agrees with the anchored games. In
this run, chained Elo misled mainly because the per-generation matches were too cheap.
We describe the laptop-scale training pipeline, whose rules and search we verified to produce identical
results in Python, Swift and the browser. We also report an external calibration: the final network won
\valGnuWins{} games against GNU~Go at its weakest level. Finally, we recommend a cheap anchored protocol
for evaluating small runs.
\end{abstract}
